\partframe{五}{Sen 的试金石与 Seiberg--Witten}

\begin{frame}{1990 年代初：弦论里的“电—磁对偶”}
  \srcp{29} 到 1990 年代初，关于 Montonen--Olive 对偶的图景在改变。主要的变化，对我来说，
  是同事们——Ashoke Sen、John Schwarz、Mike Duff、Paul Townsend、Chris Hull，以及
  不幸还有太多其他人，没法一一列举——正在为弦论（而不是量子场论）发展 Montonen--Olive
  对偶的类比。
  \vspace{0.4em}

  这些论证很有启发性。但对我来说，当时还没有“冒烟的枪”，我也看不出怎么可能取得
  扎实的进展。不过我并不像 1977 年那样怀疑了。
  \begin{beizhu}
  \textbf{批注}\quad “smoking gun”（冒烟的枪）是这场演讲的关键词：一个不依赖对偶假设、
  却能检验对偶预言的独立计算。Witten 要的不是更多证据，而是判决性的检验。
  \end{beizhu}
\end{frame}

\begin{frame}{1994 年 2 月 7 日}
  \srcp{30} 冒烟的枪，1994 年 2 月 7 日提交到 arXiv：
  \vspace{0.2em}

  \begin{tcolorbox}[enhanced,colback=black!3,colframe=black!3,boxrule=0pt,left=7pt,right=7pt,top=5pt,bottom=5pt,fontupper=\small]
  \textbf{A. Sen}, \emph{Dyon--monopole bound states, self-dual harmonic forms on the
  multi-monopole moduli space, and $SL(2,\ZZ)$ invariance in string theory},
  Physics Letters B329 (1994) 217--221.
  \vspace{0.4em}

  \emph{Existence of $SL(2,\ZZ)$ duality in toroidally compactified heterotic string
  theory (or in the $\cN=4$ supersymmetric gauge theories), that includes the strong-weak
  coupling duality transformation, implies the existence of certain supersymmetric bound
  states of monopoles and dyons. We show that the existence of these bound states, in turn,
  requires the existence of certain normalizable, (anti-)self-dual, harmonic forms on the
  moduli space of BPS multi-monopole configurations, with specific symmetry properties.
  We give an explicit construction of this harmonic form on the two monopole moduli space,
  thereby proving the existence of all the required bound states in the two-monopole sector.}
  \end{tcolorbox}
\end{frame}

\begin{frame}{Sen 的论证：三步}
  \srcp{31} Sen 明白：$\cN=4$ 超杨--米尔斯理论的 Montonen--Olive 对偶，蕴涵着多单极子
  束缚态的存在。
  \begin{enumerate}
    \item 对偶把电与磁互换，于是把 $(m,n)$ 电荷的态映到 $(n,-m)$（或类似）的态；
      $SL(2,\ZZ)$ 轨道上必须存在相应的 BPS 态。
    \item 在单极子模空间上看，这些束缚态就是模空间上中间维数的 $L^2$ 调和形式。
    \item 对两单极子，Atiyah--Hitchin 已显式描述了四维超 Kähler 模空间；它的度量足够简单，
      使 Sen 能显式地构造所需的调和形式。
  \end{enumerate}
  \begin{beizhu}
  \textbf{批注}\quad 这一步的漂亮之处：对偶预言的是“存在某些态”，而 Sen 把它翻译成
  “某个四维流形上存在某个 $L^2$ 调和 2-形式”——一个纯几何、可逐项计算的问题。
  \end{beizhu}
\end{frame}

\begin{frame}{为什么这是判决性的}
  \srcp{32} 这篇论文像电击一样击中了我。
  \begin{itemize}
    \item 一方面，我把它看作自早年（1977 年）以来对 Montonen--Olive 对偶的第一次根本性的
      新检验。在它通过这项检验之后，我真的不再认真怀疑它是否成立了。
    \item 另一方面，Sen 表明了手头的工具真的足以取得实质进展。这激励我们所有人更加努力。
  \end{itemize}
  \vspace{0.5em}
  \begin{beizhu}
  \textbf{批注}\quad 注意 Witten 的措辞：“不再认真怀疑”。他没有说“被证明”——$\cN=4$ 的
  $SL(2,\ZZ)$ 对偶至今没有数学意义上的证明；被证明的是“对偶 $\Rightarrow$ 这些态存在”，
  以及这些态确实存在。
  \end{beizhu}
\end{frame}

\begin{frame}{与 Seiberg 的五个多月}
  \srcp{33} 大约在此前一年，Nati Seiberg 在 $\cN=1$ 超杨--米尔斯理论的“动力学”
  （量子行为）上取得了重大进展。$\cN=1$ 超杨--米尔斯比 $\cN=2$ 更接近物理，但离
  Montonen--Olive 对偶更远。Nati 邀我参与进来，尝试理解 $\cN=2$ 的量子行为。
  \vspace{0.4em}

  要记得，$\cN=2$ 正是对 Montonen--Olive 对偶“几乎”正确的那个理论。Sen 的论文以及
  某些定性考虑，把我们的注意力集中到电—磁对偶在 $\cN=2$ 动力学中的作用。
  大约五个月后，我们发现：$\cN=2$ 的动力学由 Montonen--Olive 对偶的一个
  \emph{有效场论版本} 支配。
\end{frame}

\begin{frame}{$\cN=2$ 的低能有效理论}
  \srcp{34} 主要的洞见是：$\cN=2$ 理论一般可以描述为一个自由的 $U(1)$ 规范理论，
  带一些有质量的带电粒子与单极子；但在场空间的特殊点上，单极子会变得很轻，甚至无质量。
  \vspace{0.4em}

  在弱耦合下，单极子是延展对象（Bogomolny 解）；在强耦合下，它可以表现得像点状的基本粒子。
  这正是 Montonen--Olive 对偶的“有效场论”版本：
  \[
    \text{弱耦合的磁单极子}\;\longleftrightarrow\;\text{强耦合的基本粒子}.
  \]
  \begin{beizhu}
  \textbf{批注}\quad 这里发生了视角的转换：不再问“两个理论是否同一个”，而问
  “在同一个理论的某个真空附近，低能自由度是谁”。对偶变成了关于\emph{谱}的陈述。
  \end{beizhu}
\end{frame}

\begin{frame}{Seiberg--Witten 解：一条椭圆曲线}
  $\cN=2$ 纯超杨--米尔斯（规范群 $SU(2)$）的低能有效理论是 $U(1)$ 的；它的 Coulomb
  分支由一个复参数 $u=\langle\Tr\phi^2\rangle$ 参数化。Seiberg 与 Witten 发现，有效耦合
  $\tau(u)$ 不是微扰算出来的，而是某个椭圆曲线族的周期：
  \[
    y^2=x^3-u\,x^2+\Lambda^4 x ,
  \]
  其中 $\Lambda$ 是强耦合标度。$\tau(u)$ 由曲线上的全纯微分 $\dd x/y$ 的周期比给出。
  \begin{beizhu}
  \textbf{批注}\quad 曲线的判别式在 $u=\pm2\Lambda^2$ 处为零：两个奇异点。在一个点上，
  一个磁单极子变得无质量；在另一个点上，一个双子（dyon）变得无质量。$\tau(u)$ 绕奇异点
  的单值变换，正是 $SL(2,\ZZ)$ 里的 $S$ 与 $T$。
  \end{beizhu}
\end{frame}

\begin{frame}{单极子凝聚与夸克禁闭}
  在磁单极子凝聚的真空里，电场被排斥到一根弦上（对偶迈斯纳效应）：
  \[
    \text{磁单极子凝聚}\;\Longrightarrow\;\text{色通量管}\;\Longrightarrow\;
    \text{禁闭}.
  \]
  这正是 't Hooft 在 1970 年代设想的禁闭机制。这也是“强耦合下单极子可以像基本粒子”
  这句话的具体含义：在弱耦合下它是孤子，在强耦合下它是有效理论里的基本场。
  \begin{beizhu}
  \textbf{批注}\quad Seiberg--Witten 因此把两个“猜想”接上了：Montonen--Olive 的
  电—磁对偶，与 't Hooft 的“单极子凝聚导致禁闭”。在 $\cN=2$ 里，后者第一次被显式算出来。
  \end{beizhu}
\end{frame}

\begin{frame}{Seiberg--Witten 方程}
  \srcp{35} 把这些结果用到与 Donaldson 理论相关的扭曲 $\cN=2$ 上，我们得到：
  Donaldson 不变量所包含的四维流形信息，可以等价地由计数 Seiberg--Witten 方程的解得到。
  \[
    F^+=\overline{M}\,\Gamma M,\qquad \not{D}M=0 .
  \]
  这里 $A$ 是 $U(1)$ 联络，$M$ 是带电旋量场，$\Gamma$ 是 Clifford 乘法。
  \begin{beizhu}
  \textbf{批注}\quad $F^+$ 是 $A$ 的自对偶部分；第一条方程说“曲率由旋量的平方给出”，
  第二条说“旋量在联络下是调和的”。这是一阶椭圆系统，模空间有限维紧致，
  比 ASD 瞬子模空间好处理得多。
  \end{beizhu}
\end{frame}

\begin{frame}{从 Donaldson 到 Seiberg--Witten：数学的兑现}
  \begin{itemize}
    \item Donaldson 理论要数 ASD 联络模空间（维数高、奇点多）；
    \item Seiberg--Witten 只要数 $U(1)$ 联络与旋量场的解（维数为零，本质是有限个点）；
    \item 于是 Donaldson 不变量可以由 Seiberg--Witten 不变量算出。
  \end{itemize}
  \vspace{0.3em}
  这一等价后来在数学上被证明（在适当技术条件下，由 Feehan--Leness 等人完成）。
  这是本演讲所述物理图像最直接的一次“数学兑现”。
  \begin{beizhu}
  \textbf{批注}\quad 这里的逻辑值得记住：物理给出的是“两个计数问题相等”，
  数学给出的是证明。Witten 1994 年的两篇论文（\emph{Monopoles and four-manifolds}）
  把这一图像带进了四维拓扑。
  \end{beizhu}
\end{frame}

\begin{frame}{但 Langlands 呢？}
  \srcp{36} 到 1990 年代中期，已经很清楚电—磁对偶极其重要。但它与 Langlands 对应的
  关系仍然不明。
  \vspace{0.5em}

  演讲在这里转了一个弯：Witten 回到另一条线——几何 Langlands 对应。
  \begin{center}
  \begin{tikzpicture}[>={Stealth[length=2mm]}]
    \node[draw,rounded corners,inner sep=6pt] (a) at (0,0) {电磁对偶};
    \node[draw,rounded corners,inner sep=6pt] (b) at (4.6,0) {超对称 / BPS};
    \node[draw,rounded corners,inner sep=6pt] (c) at (9.2,0) {几何 Langlands};
    \draw[->] (a) -- (b);
    \draw[->] (b) -- (c);
    \draw[->,dashed] (a) to[bend right=28] node[below,font=\scriptsize]{1977 年 Atiyah 的那句话} (c);
  \end{tikzpicture}
  \end{center}
\end{frame}

\partframe{六}{几何 Langlands 与 Hecke 变换}

\begin{frame}{Beilinson--Drinfeld 的“几何”版本}
  \srcp{37} 在继续之前，我得回溯并解释另一条线索。大约从 1990 年起，Alexander Beilinson
  与 Vladimir Drinfeld 对普通 Riemann 面（而不是数域）表述了 Langlands 对应的
  “几何”版本。
  \vspace{0.4em}

  从一开始，他们就大量使用物理学里熟悉的东西，例如流代数（仿射 Kac--Moody 代数）。
  我费了很大力气去理解他们在说什么，但在那些年里我觉得这不可能。物理里熟悉的东西，
  似乎以一种完全不熟悉的方式被搅乱了。
  \begin{beizhu}
  \textbf{批注}\quad 数域情形：Galois 表示 $\leftrightarrow$ 自守表示。
  几何情形：$C$ 上的 $\gvee$-局部系统 $\leftrightarrow$ $\Bun_G$ 上的 D-模。
  “几何”指的是把函数换成层、把特征值换成 Hecke 特征层。
  \end{beizhu}
\end{frame}

\begin{frame}{Hecke 修改}
  \srcp{38} 设 $E\to C$ 是 $C$ 上的全纯 $G_\CC$-丛（$G_\CC$ 是单复群）。
  $E$ 在点 $p\in C$ 处的 \emph{Hecke 修改}，是另一个全纯 $G_\CC$-丛 $E'\to C$，
  连同 $C\setminus p$ 上的一个同构
  \[
    \varphi:\;E|_{C\setminus p}\;\xrightarrow{\ \sim\ }\;E'|_{C\setminus p}.
  \]
  \begin{center}
  \begin{tikzpicture}[scale=0.95,>={Stealth[length=1.5mm]}]
    \draw[thick] (0,0) .. controls (2,0.9) and (4.5,-0.9) .. (6.4,0.1);
    \fill (3.2,0.05) circle (1.7pt) node[below=1pt,font=\scriptsize]{$p$};
    \node[font=\scriptsize] at (0.5,0.55) {$C$};
    \node[font=\scriptsize,align=left] at (7.9,0.75) {$E$};
    \node[font=\scriptsize,align=left] at (7.9,-0.55) {$E'$};
    \node[font=\scriptsize] at (3.2,-0.95) {$\varphi:E\cong E'$ 于 $C\setminus p$};
  \end{tikzpicture}
  \end{center}
  这就是几何 Langlands 的“基本操作”：除了一个点，两个丛完全一样。
\end{frame}

\begin{frame}{两个例子}
  \textbf{例 1：$G_\CC=\CC^*$}。全纯线丛 $E=\cL\to C$，在 $p$ 处的修改是
  \[
    \cL'=\cL(p)=\cL\otimes\cO(p),
  \]
  按定义，它只在 $p$ 处与 $\cL$ 不同。
  \vspace{0.6em}

  \textbf{例 2：$G_\CC=SL(2,\CC)$}。$E$ 是秩 $2$ 全纯向量丛。
  若局部写成两个线丛之和 $E=\cL_1\oplus\cL_2$，则一个 Hecke 修改是
  \[
    E'=\cL_1(p)\oplus\cL_2(-p).
  \]
  \begin{beizhu}
  \textbf{批注}\quad 两个例子的共同点：修改只改变 $p$ 处的“局部格”，
  而 $\cL_1(p)\oplus\cL_2(-p)$ 与 $\cL_1\oplus\cL_2$ 在 $C\setminus p$ 上典范同构。
  不同的修改方式，对应不同的“型”。
  \end{beizhu}
\end{frame}

\begin{frame}{型：Hecke 修改的分类}
  \srcp{39} 在给定点 $p$，Hecke 修改可以有不同的类型。例如对 $G_\CC=\CC^*$，
  除了 $\cL\mapsto\cL(p)$，还可以取 $\cL(np)=\cL\otimes\cO(p)^n$，$n\in\ZZ$ 任意。
  \vspace{0.4em}

  一般的 $G$：Hecke 修改的型由\emph{支配余权} $\lambda\in X_*(T)$（模 Weyl 群）给出：
  \[
    \{\text{Hecke 修改的型}\}\;\cong\;X_*(T)/W .
  \]
  \begin{beizhu}
  \textbf{批注}\quad 几何上：$E$ 在 $p$ 附近平凡化后，修改就是把平凡丛的局部截面格
  $\cO^{\oplus r}$ 换成另一个格 $\cL_\lambda$；型的集合是仿射 Grassmannian 的轨道。
  \end{beizhu}
\end{frame}

\begin{frame}{几何 Satake：型由对偶群的表示标号}
  \srcp{39} 几何 Langlands 纲领的一个重要输入，是\emph{几何 Satake 对应}
  （由 Lusztig、Drinfeld、Ginzburg、Mirković、Vilonen 等人发展）：
  \[
    \Perv_{G(\cO)}(\Gr_G)\;\simeq\;\Rep(\gvee),
  \]
  其中 $\Gr_G=G(\CC((t)))/G(\CC[[t]])$ 是仿射 Grassmannian，$\cO=\CC[[t]]$。
  \vspace{0.4em}

  换句话说：Hecke 修改的类型，由对偶群 $\gvee$ 的不可约有限维表示来标号。
  对 $G_\CC=\CC^*$，对偶群仍是 $\CC^*$，不可约表示由整数 $n$ 标号——这正是
  $\cL\mapsto\cL(np)$ 里的 $n$。
  \begin{beizhu}
  \textbf{批注}\quad 这是 $\gvee$ 第一次以“表示论”的方式进入：不是通过磁荷，
  而是通过 Hecke 修改的类型。两种进入方式在 1977 年是同一件事，在这里又重逢了。
  \end{beizhu}
\end{frame}

\begin{frame}{Hecke 特征层}
  设 $\cF$ 是 $\Bun_G$ 上的一个（holonomic）D-模。对每个点 $x\in C$ 与 $\gvee$ 的
  不可约表示 $V$，Hecke 对应给出
  \[
    H^V_x(\cF)\;\in\;D\text{-}\mathrm{mod}(\Bun_G),
  \]
  若对所有 $x,V$ 都有
  \[
    H^V_x(\cF)\;\cong\;V\otimes\cF ,
    \qquad (V\otimes\cF:\;\text{把 }V\text{ 看成常系数}) ,
  \]
  则称 $\cF$ 是\emph{Hecke 特征层}。
  \vspace{0.4em}

  几何 Langlands 猜想：对 $C$ 上每个 $\gvee$-局部系统 $\cE$，存在这样的 $\cF_\cE$，
  且其“特征值”就是 $\cE$。这就是“自守 $\leftrightarrow$ Galois”的几何版。
\end{frame}

\begin{frame}{多年困惑：Hecke 与物理何干？}
  \srcp{40} 多年来，我一直困惑：这样的几何 Hecke 变换怎么可能与物理有关？
  我假定它们必须有关——如果几何 Langlands 要有物理解释的话。
  \vspace{0.5em}

  答案最终出奇地简单，用到了量子规范理论的基本构件：
  \begin{itemize}
    \item Wilson 与 't Hooft 算子的电—磁对偶；
    \item Atiyah 与 Hitchin 等人彻底研究过的 Bogomolny 方程。
  \end{itemize}
  这是 Witten 2006 年与 Anton Kapustin 合写的“规范理论与几何 Langlands”一文的主要想法之一。
\end{frame}

\begin{frame}{本章小结：三件事在同一个点上重合}
  \begin{center}
  \renewcommand{\arraystretch}{1.3}
  \begin{tabular}{@{}lll@{}}
    \toprule
    物理 & 群论 & 几何 \\
    \midrule
    电电荷 & $G$ 的表示 $R$ & $X^*(T)$ 里的最高权 \\
    磁荷 & $\gvee$ 的表示 & $X_*(T)$ 里的余权 \\
    Hecke 修改的型 & $\gvee$ 的不可约表示 & 支配余权 $\lambda$ \\
    \bottomrule
  \end{tabular}
  \end{center}
  \vspace{0.5em}
  接下来的三部分，就是把这张表一行一行地变成定理。
\end{frame}
